% Source: Weinberg GR, physical PDF pp. 573--589
%         (printed pp. 545--561; boundary pages are partial).
% Coverage: Section 15.7, Helium Synthesis; equations (15.7.1)--(15.7.37),
%           all intervening unnumbered displays, and Tables 15.5--15.7.
% Figures/tables/footnotes: Tables 15.5--15.7 with source notes; linked
%           reference markers 52, 93, 99--142; no figures or footnotes.
% Status: source-reviewed and compile-clean; visually proofed.
% Uncertainties: the unit and differential-measure modernizations are
%                documented below; no unresolved source uncertainty remains.
%
% MODERNIZATION CHECK: units are set upright and differential measures use
% \dd. Source-era reaction and isotope notation is otherwise retained.
% No signature-, coordinate-order-, scale-factor-, action-, or worldline-
% derivative-dependent expression occurs in this section.

\subsection{Helium Synthesis}\label{sec:15.7}

The relative abundance of the chemical elements has been under careful study
by geologists and astronomers since the pioneering work of Frank Wigglesworth
Clarke\hyperref[ch15-ref-99]{\textsuperscript{99}} in the last century.
Gradually these studies have revealed a ``cosmic'' distribution of
abundances,\hyperref[ch15-ref-100]{\textsuperscript{100}} with hydrogen and
then helium by far the most abundant elements, followed by the group
C--N--O--Ne, and with the group Li--Be--B and all elements heavier than nickel
scarce. The problem of explaining these abundances has long been considered
one of the major challenges facing theoretical astrophysics.

One possible explanation lay in the nuclear reactions that provide energy to
the stars. Rutherford's demonstration of nuclear transmutation in the
laboratory led Eddington\hyperref[ch15-ref-101]{\textsuperscript{101}} in
1920 to suggest that the sun might derive its energy from the fusion of
hydrogen into helium. If so, then perhaps the stars (or at least the first
generation of stars) were formed from pure hydrogen, and have gradually
produced helium and heavier elements as ashes of their internal fires. The
detailed reactions by which the stars burn hydrogen to helium were laid out
in 1939 by Hans Bethe,\hyperref[ch15-ref-102]{\textsuperscript{102}} and the
subsequent reactions in which helium fuses into heavier elements were explored
in the 1950's in a series of papers by
Salpeter,\hyperref[ch15-ref-103]{\textsuperscript{103}} the Burbidges,
Fowler, and Hoyle,\hyperref[ch15-ref-104]{\textsuperscript{104}}
Cameron,\hyperref[ch15-ref-105]{\textsuperscript{105}} and others. Most
recently, Clayton and Arnett\hyperref[ch15-ref-106]{\textsuperscript{106}}
have emphasized the importance of stellar explosions as an agent in
nucleosynthesis.

There is one other competing theory of nucleosynthesis, worked out in the
late 1940's by George Gamow and his
collaborators.\hyperref[ch15-ref-107]{\textsuperscript{107}} Gamow reasoned
that, although the early hot dense period of cosmic expansion was much briefer
than the lifetime of a star, there was a large number of free neutrons present
at that time, so that the heavy elements could be built up quickly by
successive neutron captures, starting with
\(n+p\to d+\gamma\). The abundances of the elements would then be correlated
with their neutron capture cross-sections, in rough agreement with
observation. We have already noted in Section~\ref{sec:15.5} that the
necessity of avoiding too much helium production in this theory required the
presence of black-body radiation, with a present temperature that was
estimated\hyperref[ch15-ref-52]{\textsuperscript{52}} as
\(5^\circ\mathrm K\).

Both the stellar and the cosmological theories of nucleosynthesis have their
limitations. There are no stable nuclei with atomic weights \(A=5\) or
\(A=8\), so it is difficult to build up elements heavier than helium by
\(p\)-\(\alpha\), \(n\)-\(\alpha\), or \(\alpha\)-\(\alpha\) collisions.
In stars that have converted all hydrogen to helium at their cores, it is
possible to bridge the gaps at \(A=5\) and \(A=8\) by the production of small
amounts of the unstable nuclide \(\mathrm{Be}^8\) in
\(\alpha\)-\(\alpha\) collisions, followed by the production of
\(\mathrm C^{12}\) in \(\alpha\)-\(\mathrm{Be}^8\) collisions.%
\hyperref[ch15-ref-103]{\textsuperscript{103}} However, the density of the
expanding universe at the temperature \(T\simeq10^{9\,\circ}\mathrm K\) is
too low to allow much helium burning to occur. It is generally accepted today
that all elements heavier than helium were synthesized in stars.

On the other hand, several
authors\hyperref[ch15-ref-108]{\textsuperscript{108}} have noted that the
cosmic abundance of helium is too large to be easily explained in terms of
nucleosynthesis in stars. The luminosity-to-mass ratio \(L/M\) of our galaxy
is about one-tenth the solar ratio \(L_\odot/M_\odot\), or
\(0.2\,\mathrm{erg/(g\,sec)}\). If the luminosity of the galaxy has remained
constant during the last \(10^{10}\) years, then about
\(0.06\,\mathrm{MeV}\) per nucleon would have been produced. In contrast,
the fusion of hydrogen into helium releases about \(6\,\mathrm{MeV}\) per
nucleon, so not more than about 1\% of the nucleons in our galaxy could have
been fused into helium (or heavier nuclei) by ordinary stellar processes. As
we shall see, estimates of the present helium abundance vary, but there is
wide agreement that the cosmic abundance of helium by mass is considerably
greater than 1\%. It is of course possible that the helium could have been
synthesized in an earlier, more luminous, epoch of our galaxy; as already
remarked in Section~\ref{sec:15.5}, the released energy could, if thermalized,
account for the present \(2.7^\circ\mathrm K\) microwave background. However,
it is more interesting and more natural to assume that the large cosmic
helium abundance was produced during the early history of the universe, with
the energy of fusion mostly lost in the subsequent red shift.

Let us now calculate the cosmologically produced abundance of helium. It is
very convenient to divide this calculation into two parts. First, we calculate
the neutron--proton abundance ratio as a function of time, taking account only
of the weak interaction processes
\begin{equation}
  n+\nu\rightleftarrows p+e^-,
  \qquad
  n+e^+\rightleftarrows p+\bar\nu,
  \qquad
  n\rightleftarrows p+e^-+\bar\nu.
  \tag{15.7.1}\label{eq:15.7.1}
\end{equation}
(Here \(\nu\) will mean \(\nu_e\).) In the second part of this calculation,
we put in the nuclear reactions that lead to helium synthesis.

The number densities of \(\nu,\bar\nu,e^-\), and \(e^+\) are given here by
the Fermi distributions \eqref{eq:15.6.3}, with zero chemical potential and
with different temperatures \(T\) or \(T_\nu\) for \(e^\pm\) (and \(\gamma\))
or \(\nu\) and \(\bar\nu\):
\[
  \begin{aligned}
    n_{e^-}(p)\dd p
    &=
    n_{e^+}(p)\dd p
    =
    8\pi h^{-3}p^2\dd p
    \left[
      \exp\left(\frac{E_e(p)}{kT}\right)+1
    \right]^{-1},
    \\
    n_\nu(p)\dd p
    &=
    n_{\bar\nu}(p)\dd p
    =
    4\pi h^{-3}p^2\dd p
    \left[
      \exp\left(\frac{E_\nu(p)}{kT_\nu}\right)+1
    \right]^{-1},
  \end{aligned}
\]
where
\[
  E_e(p)=(p^2+m_e^2)^{1/2},
  \qquad
  E_\nu(p)=p.
\]

The rates for the various reactions \eqref{eq:15.7.1} are given by the
``\(V-A\)'' theory of weak
interactions,\hyperref[ch15-ref-109]{\textsuperscript{109}} except that the
Pauli exclusion principle suppresses these rates by a factor equal to the
fraction of all states that are unfilled:
\[
  \begin{aligned}
    1-
    \left[
      \exp\left(\frac{E_e}{kT}\right)+1
    \right]^{-1}
    &=
    \left[
      1+\exp\left(-\frac{E_e}{kT}\right)
    \right]^{-1},
    \\
    1-
    \left[
      \exp\left(\frac{E_\nu}{kT_\nu}\right)+1
    \right]^{-1}
    &=
    \left[
      1+\exp\left(-\frac{E_\nu}{kT_\nu}\right)
    \right]^{-1}.
  \end{aligned}
\]
The rates (per nucleon) of the processes \eqref{eq:15.7.1} are then
\begin{align}
  \lambda(n+\nu\to p+e^-)
  &=
  A\int
  v_eE_e^2p_\nu^2\dd p_\nu
  [e^{E_\nu/kT_\nu}+1]^{-1}
  [1+e^{-E_e/kT}]^{-1},
  \tag{15.7.2}\label{eq:15.7.2}\\
  \lambda(n+e^+\to p+\bar\nu)
  &=
  A\int
  E_\nu^2p_e^2\dd p_e
  [e^{E_e/kT}+1]^{-1}
  [1+e^{-E_\nu/kT_\nu}]^{-1},
  \tag{15.7.3}\label{eq:15.7.3}\\
  \lambda(n\to p+e^-+\bar\nu)
  &=
  A\int
  v_eE_\nu^2E_e^2\dd p_\nu
  [1+e^{-E_\nu/kT_\nu}]^{-1}
  [1+e^{-E_e/kT}]^{-1},
  \tag{15.7.4}\label{eq:15.7.4}\\
  \lambda(p+e^-\to n+\nu)
  &=
  A\int
  E_\nu^2p_e^2\dd p_e
  [e^{E_e/kT}+1]^{-1}
  [1+e^{-E_\nu/kT_\nu}]^{-1},
  \tag{15.7.5}\label{eq:15.7.5}\\
  \lambda(p+\bar\nu\to n+e^+)
  &=
  A\int
  v_eE_e^2p_\nu^2\dd p_\nu
  [e^{E_\nu/kT_\nu}+1]^{-1}
  [1+e^{-E_e/kT}]^{-1},
  \tag{15.7.6}\label{eq:15.7.6}\\
  \lambda(p+e^-+\bar\nu\to n)
  &=
  A\int
  v_eE_e^2p_\nu^2\dd p_\nu
  [e^{E_e/kT}+1]^{-1}
  [e^{E_\nu/kT_\nu}+1]^{-1}.
  \tag{15.7.7}\label{eq:15.7.7}
\end{align}
Here \(A\) is the constant
\begin{equation}
  A=\frac{g_V^2+3g_A^2}{2\pi^3\hbar^7},
  \tag{15.7.8}\label{eq:15.7.8}
\end{equation}
where \(g_V\) and \(g_A\) are the vector and axial-vector coupling constants
of the nucleon, taken here to have the values
\begin{equation}
  g_V=1.418\times10^{-49}\,\mathrm{erg\,cm^3},
  \qquad
  g_A=1.18g_V.
  \tag{15.7.9}\label{eq:15.7.9}
\end{equation}

Also, \(E_e\) and \(E_\nu\) are related by
\begin{align}
  E_e-E_\nu&=Q
  &&\text{for \(n+\nu\rightleftarrows p+e^-\),}
  \tag{15.7.10}\label{eq:15.7.10}\\
  E_\nu-E_e&=Q
  &&\text{for \(n+e^+\rightleftarrows p+\bar\nu\),}
  \tag{15.7.11}\label{eq:15.7.11}\\
  E_\nu+E_e&=Q
  &&\text{for \(n\rightleftarrows p+e^-+\bar\nu\),}
  \tag{15.7.12}\label{eq:15.7.12}
\end{align}
where
\begin{equation}
  Q\equiv m_n-m_p=1.293\,\mathrm{MeV}.
  \tag{15.7.13}\label{eq:15.7.13}
\end{equation}
The integrals \eqref{eq:15.7.2}--\eqref{eq:15.7.7} are taken over the
positive values of \(p_\nu\) and \(p_e\) allowed by these relations. It is
very convenient to write all integrals in terms of a variable of integration
\(q\), taken as \(E_\nu\) in Eqs.~\eqref{eq:15.7.2},
\eqref{eq:15.7.4}, and \eqref{eq:15.7.5}, and as \(-E_\nu\) in
Eqs.~\eqref{eq:15.7.3}, \eqref{eq:15.7.6}, and \eqref{eq:15.7.7}.
Replacing \(p_e^2\dd p_e\) with \(v_eE_e^2\dd E_e\), the total
\(n\to p\) and \(p\to n\) transition rates are then
\begin{align}
  \lambda(n\to p)
  &=
  \lambda(n+\nu\to p+e^-)
  +\lambda(n+e^+\to p+\bar\nu)
  +\lambda(n\to p+e^-+\bar\nu)
  \notag\\
  &=
  A\int
  \left[
    1-\frac{m_e^2}{(Q+q)^2}
  \right]^{1/2}
  (Q+q)^2q^2\dd q
  [1+e^{q/kT_\nu}]^{-1}
  [1+e^{-(Q+q)/kT}]^{-1},
  \tag{15.7.14}\label{eq:15.7.14}\\
  \lambda(p\to n)
  &=
  \lambda(p+e^-\to n+\nu)
  +\lambda(p+\bar\nu\to n+e^+)
  +\lambda(p+e^-+\bar\nu\to n)
  \notag\\
  &=
  A\int
  \left[
    1-\frac{m_e^2}{(Q+q)^2}
  \right]^{1/2}
  (Q+q)^2q^2\dd q
  [1+e^{-q/kT_\nu}]^{-1}
  [1+e^{(Q+q)/kT}]^{-1}.
  \tag{15.7.15}\label{eq:15.7.15}
\end{align}
The integrals now run from \(-\infty\) to \(+\infty\), leaving out a gap
from \(-Q-m_e\) to \(-Q+m_e\). The differential equation for the ratio
\(X_n\) of neutrons to all nucleons is
\begin{equation}
  -\frac{\dd X_n}{\dd t}
  =
  \lambda(n\to p)X_n
  -
  \lambda(p\to n)(1-X_n).
  \tag{15.7.16}\label{eq:15.7.16}
\end{equation}
The solution of this equation has been calculated by
Peebles,\hyperref[ch15-ref-93]{\textsuperscript{93}} and is presented here
in Table~\ref{tab:15.5}. Although the quantitative behavior of \(X_n(t)\)
can only be obtained by a numerical integration, it is possible to appreciate
the main features of this solution through a few qualitative observations:

\emph{(A)} For \(kT\gg Q\), we can set \(T=T_\nu\) and put \(Q\) and \(m_e\)
equal to zero in Eqs.~\eqref{eq:15.7.14} and \eqref{eq:15.7.15}. The
transition rates are then
\begin{align}
  \lambda(n\to p)
  \simeq
  \lambda(p\to n)
  &\simeq
  A\int_{-\infty}^{\infty}
  q^4\dd q\,
  (1+e^{-q/kT})^{-1}
  (1+e^{q/kT})^{-1}
  \notag\\
  &=
  \frac{7}{15}\pi^4A(kT)^5
  =
  0.361\,\mathrm{sec}^{-1}
  \left(
    \frac{T}{10^{10\,\circ}\mathrm K}
  \right)^5.
  \tag{15.7.17}\label{eq:15.7.17}
\end{align}

\begingroup
\renewcommand{\thetable}{15.5}
\begin{table}[htbp]
  \centering
  \caption{Neutron Fraction \(X_n\), as a Function of Temperature or Time,
  with Neglect of the Formation of Complex Nuclei.}
  \label{tab:15.5}
  \small
  \begin{tabular}{@{}rrr@{}}
    \toprule
    \(T\) (\(^{\circ}\mathrm K\)) & \(t\) (sec) & \(X_n\) \\
    \midrule
    \(10^{12}\)        & \(0\)        & \(0.496\) \\
    \(3\times10^{11}\) & \(0.001129\) & \(0.488\) \\
    \(10^{11}\)        & \(0.01078\)  & \(0.462\) \\
    \(3\times10^{10}\) & \(0.1209\)   & \(0.380\) \\
    \(10^{10}\)        & \(1.103\)    & \(0.241\) \\
    \(3\times10^9\)    & \(13.83\)    & \(0.170\) \\
    \(1.3\times10^9\)  & \(98^*\)     & \(0.150\) \\
    \(1.2\times10^9\)  & \(119^*\)    & \(0.147\) \\
    \(1.1\times10^9\)  & \(146^*\)    & \(0.143\) \\
    \(1.0\times10^9\)  & \(182.0\)    & \(0.137\) \\
    \(9\times10^8\)    & \(226^*\)    & \(0.131\) \\
    \(8\times10^8\)    & \(290^*\)    & \(0.123\) \\
    \(7\times10^8\)    & \(383^*\)    & \(0.112\) \\
    \(3\times10^8\)    & \(2080\)     & \(0.021\) \\
    \(10^8\)           & \(18700\)    & \(10^{-8}\) \\
    \bottomrule
  \end{tabular}

  \vspace{0.5em}
  \begin{minipage}{0.92\linewidth}
    \footnotesize
    Values of \(t\) are taken from the calculations of P.~J.~E. Peebles,
    \emph{Astron.\ J.} \textbf{146}, 542 (1966), except that values marked
    with an asterisk are interpolated from Peebles' results. Values of
    \(X_n\), for \(T>1.0\times10^{9\,\circ}\mathrm K\), are taken from
    Peebles, op.\ cit., Table 4. Values of \(X_n\), for
    \(T<10^{9\,\circ}\mathrm K\), are calculated from the value at
    \(10^{9\,\circ}\mathrm K\), under the assumption that \(X_n\) decreases
    exponentially at the rate \((1013\,\mathrm{sec})^{-1}\) of free neutron
    decay.
  \end{minipage}
\end{table}
\endgroup

This may be compared with the ``age'' \(t\), given by
Eqs.~\eqref{eq:15.6.44} and \eqref{eq:15.6.33}:
\begin{equation}
  t
  =
  1.09\,\mathrm{sec}
  \left(
    \frac{T}{10^{10\,\circ}\mathrm K}
  \right)^{-2}.
  \tag{15.7.18}\label{eq:15.7.18}
\end{equation}
We see that the product \(\lambda t\) is larger than 10 for
\(T\gtrsim3\times10^{10\,\circ}\mathrm K\), so at these temperatures the
neutron fraction \(X_n\) should be given by the equilibrium solution of
Eq.~\eqref{eq:15.7.16}, which is
\begin{equation}
  X_n
  =
  \frac{\lambda(p\to n)}
       {\lambda(p\to n)+\lambda(n\to p)}.
  \tag{15.7.19}\label{eq:15.7.19}
\end{equation}
Note that Eq.~\eqref{eq:15.7.17} will not be quantitatively correct when
\(T\) drops to near \(3\times10^{10\,\circ}\mathrm K\), because \(kT\) is
then not very much larger than \(Q\). However, even though the rates
\(\lambda(p\to n)\) and \(\lambda(n\to p)\) may differ somewhat from
Eq.~\eqref{eq:15.7.17} and from each other, they are still large enough for
\(T\gtrsim3\times10^{10\,\circ}\mathrm K\) to justify the use of the
equilibrium solution \eqref{eq:15.7.19}.

\emph{(B)} As long as \(T_\nu\simeq T\) (that is, for
\(T>10^{10\,\circ}\mathrm K\)) the rates \eqref{eq:15.7.14} and
\eqref{eq:15.7.15} have the ratio
\begin{equation}
  \frac{\lambda(p\to n)}{\lambda(n\to p)}
  \simeq
  \exp\left(-\frac{Q}{kT}\right).
  \tag{15.7.20}\label{eq:15.7.20}
\end{equation}
Thus Eq.~\eqref{eq:15.7.19} gives the neutron abundance for
\(T\gtrsim3\times10^{10\,\circ}\mathrm K\) as
\begin{equation}
  X_n\simeq[1+e^{Q/kT}]^{-1}.
  \tag{15.7.21}\label{eq:15.7.21}
\end{equation}
The neutron abundance starts at \(X_n\simeq\tfrac12\) at very early times,
and drops slowly as the temperature falls, reaching \(X_n\simeq0.38\) for
\(T=3\times10^{10\,\circ}\mathrm K\). It is a profoundly important fact
that the initial condition for Eq.~\eqref{eq:15.7.16} does not have to be
chosen arbitrarily, and does not depend on any detailed model of the very
early universe, but follows directly from the singular behavior of the rates
\(\lambda\) as \(t\to0\).%
\hyperref[ch15-ref-109a]{\textsuperscript{109a}}

\emph{(C)} Once \(T\) drops to about \(1.3\times10^{9\,\circ}\mathrm K\),
the rates of the two- and three-body reactions
\(n+\nu\rightleftarrows p+e^-\),
\(n+e^+\rightleftarrows p+\bar\nu\), and
\(p+e^-+\bar\nu\to n\) become negligible. The only remaining reaction is
the ``one-body'' process \(n\to p+e^-+\bar\nu\), which at these low
temperatures proceeds at the rate of free neutron decay, taken here to have
the value used by Peebles:\hyperref[ch15-ref-93]{\textsuperscript{93}}
\begin{equation}
  \lambda^{-1}(n\to p+e^-+\bar\nu)=1013\,\mathrm{sec}.
  \tag{15.7.22}\label{eq:15.7.22}
\end{equation}
Thus the neutron abundance from the time when
\(T\simeq1.3\times10^{9\,\circ}\mathrm K\) to the start of nucleosynthesis
is given by
\begin{equation}
  X_n(t)
  =
  N\exp\left[-\frac{t\,(\mathrm{sec})}{1013}\right].
  \tag{15.7.23}\label{eq:15.7.23}
\end{equation}
The only part of the theory of helium synthesis in which detailed numerical
calculations are really needed is the evaluation of the constant \(N\). In
carrying out this calculation, it is convenient first to ignore neutron decay
as well as nucleosynthesis, in which case the neutron abundance is a function
\(X_n^{(0)}(t)\) that approaches a finite limit as \(t\to\infty\). (This is
the quantity called \(X_n\) in Peebles' Table 1.%
\hyperref[ch15-ref-93]{\textsuperscript{93}} Peebles also ignores the
process \(p+e^-+\bar\nu\to n\), which is in fact negligible over the whole
period of interest.) Neutron decay has a negligible effect until
\(t\simeq20\,\mathrm{sec}\), whereas after that time the temperature is
below \(3\times10^{9\,\circ}\mathrm K\), so that the rate
\(\lambda(p\to n)\) is negligible compared with \(\lambda(n\to p)\), and
lepton degeneracy has little effect on the rate of neutron decay. It follows
that the whole effect of neutron decay is to multiply \(X_n^{(0)}(t)\) with
an exponential decay factor:
\begin{equation}
  X_n(t)
  \simeq
  X_n^{(0)}(t)
  \exp\left[-\frac{t\,(\mathrm{sec})}{1013}\right].
  \tag{15.7.24}\label{eq:15.7.24}
\end{equation}
Peebles\hyperref[ch15-ref-93]{\textsuperscript{93}} finds that
\(X_n^{(0)}\) approaches the value 0.1640 as \(t\to\infty\), so comparing
\eqref{eq:15.7.23} with \eqref{eq:15.7.24}, we have
\begin{equation}
  N\simeq X_n^{(0)}(\infty)=0.1640.
  \tag{15.7.25}\label{eq:15.7.25}
\end{equation}

Now we can proceed to the second part of our calculation, and put in the
nuclear reactions that lead to the synthesis of complex nuclei. At early
times, when \(T\gg10^{10\,\circ}\mathrm K\), the various nuclei would be in
thermal equilibrium, with the number density \(n_i\) of nuclide \(i\) given
by Eq.~\eqref{eq:15.6.3}. Since the nuclei are highly nonrelativistic and
nondegenerate during the whole period that concerns us here, we can use the
Maxwell--Boltzmann approximation to Eq.~\eqref{eq:15.6.3}, and write for the
total number density of species \(i\):
\begin{align}
  n_i
  &=
  4\pi g_i h^{-3}
  \exp\left(\frac{\mu_i-m_i}{kT}\right)
  \int_0^\infty q^2\dd q\,
  \exp\left(-\frac{q^2}{2m_i kT}\right)
  \notag\\
  &=
  g_i
  \left(\frac{2\pi m_i kT}{h^2}\right)^{3/2}
  \exp\left(\frac{\mu_i-m_i}{kT}\right).
  \tag{15.7.26}\label{eq:15.7.26}
\end{align}
Of course, we are not given the chemical potentials \(\mu_i\), but we know
that they are conserved in all reactions. Hence, if nuclear reactions can
rapidly build up a nucleus \(i\) out of \(Z_i\) protons and \(A_i-Z_i\)
neutrons, then \(\mu_i\) is given by
\begin{equation}
  \mu_i=Z_i\mu_p+(A_i-Z_i)\mu_n.
  \tag{15.7.27}\label{eq:15.7.27}
\end{equation}
It is convenient to write \eqref{eq:15.7.26} as a relation between the
fractions by weight of nuclide \(i\), free neutrons, and free protons:
\[
  X_i\equiv\frac{n_iA_i}{n_N},
  \qquad
  X_n\equiv\frac{n_n}{n_N},
  \qquad
  X_p\equiv\frac{n_p}{n_N},
\]
where \(n_N\) is the total number density of nucleons, bound or free:
\[
  n_N
  =
  n_{N0}\left(\frac{a_0}{a}\right)^3
  =
  \frac{\rho_{N0}}{m_N}
  \left(\frac{a_0}{a}\right)^3.
\]
Using Eq.~\eqref{eq:15.7.27}, and approximating
\(m_p=m_n=m_N\) and \(m_i=A_im_N\) in the \(3/2\)-power in
\eqref{eq:15.7.26}, we have then
\begin{equation}
  X_i
  =
  \frac12
  X_p^{Z_i}
  X_n^{A_i-Z_i}
  g_iA_i^{1/2}
  \epsilon^{A_i-1}
  \exp\left(\frac{B_i}{kT}\right),
  \tag{15.7.28}\label{eq:15.7.28}
\end{equation}
where \(B_i\) is the binding energy
\begin{equation}
  -B_i\equiv m_i-Z_im_p-(A_i-Z_i)m_n,
  \tag{15.7.29}\label{eq:15.7.29}
\end{equation}
and \(\epsilon\) is the dimensionless quantity
\begin{align}
  \epsilon
  &\equiv
  \frac12h^3n_N(2\pi m_NkT)^{-3/2}
  \notag\\
  &=
  1.61\times10^{-12}
  \left(
    \frac{\rho_{N0}}{10^{-30}\,\mathrm{g/cm^3}}
  \right)
  \left(
    \frac{a}{10^{-10}a_0}
  \right)^{-3}
  \left(
    \frac{T}{10^{9\,\circ}\mathrm K}
  \right)^{-3/2}.
  \tag{15.7.30}\label{eq:15.7.30}
\end{align}
Since \(\epsilon\) is very small in the period of interest, the abundance of
a given complex nuclide \(i\) will be very small until \(T\) drops to the
value
\begin{equation}
  T_i\simeq\frac{B_i}{k(A_i-1)\lvert\ln\epsilon\rvert}.
  \tag{15.7.31}\label{eq:15.7.31}
\end{equation}
Values of \(T_i\) are given for various nuclides and various values of the
present density \(\rho_{N0}\) in Table~\ref{tab:15.6}. Note that \(T_i\)
depends only very weakly on the present density \(\rho_{N0}\), because
\(\rho_{N0}\) enters only in the quantity \(\lvert\ln\epsilon\rvert\), and
this quantity has a value between 25 and 35 over the whole range of relevant
temperatures and densities.

\begingroup
\renewcommand{\thetable}{15.6}
\begin{table}[htbp]
  \centering
  \caption{Values for the Temperature \(T_i\) Defined by
  Eq.~\eqref{eq:15.7.31}, for Various Nuclides and Various Values of the
  Present Density \(\rho_{N0}\).}
  \label{tab:15.6}
  \small
  \setlength{\tabcolsep}{5pt}
  \begin{tabular}{@{}ccccc@{}}
    \toprule
    & \(\dfrac{B}{k(A-1)}\) &
    \multicolumn{3}{c}{\(T_i\) (\(10^{9\,\circ}\mathrm K\))} \\
    \cmidrule(lr){3-5}
    Nuclide &
    \(10^{9\,\circ}\mathrm K\) &
    \(\rho_{N0}=10^{-29}\,\mathrm{g/cm^3}\) &
    \(\rho_{N0}=10^{-30}\,\mathrm{g/cm^3}\) &
    \(\rho_{N0}=10^{-31}\,\mathrm{g/cm^3}\) \\
    \midrule
    \(\mathrm H^2\)          & 25.8 & 0.83 & 0.77 & 0.72 \\
    \(\mathrm H^3\)          & 49.3 & 1.6  & 1.5  & 1.4 \\
    \(\mathrm{He}^3\)        & 44.6 & 1.4  & 1.3  & 1.2 \\
    \(\mathrm{He}^4\), etc.  & 109  & 3.9  & 3.6  & 3.3 \\
    \bottomrule
  \end{tabular}

  \vspace{0.5em}
  \begin{minipage}{0.92\linewidth}
    \footnotesize
    Heavier nuclides have about the same values of \(T_i\) as
    \(\mathrm{He}^4\). In thermal equilibrium, \(T_i\) is the maximum
    temperature at which a nuclide \(i\) could be abundant.
  \end{minipage}
\end{table}
\endgroup

If nuclear abundances really were governed by the conditions of thermal
equilibrium down to temperatures of order \(10^{9\,\circ}\mathrm K\), then
according to Table~\ref{tab:15.6}, we should expect \(\mathrm{He}^4\) and
heavier nuclides to appear first, followed by \(\mathrm{He}^3\) and
\(\mathrm H^3\), and finally by \(\mathrm H^2\). However, this is not at all
what happens, because thermal equilibrium is not maintained down to
\(10^{9\,\circ}\mathrm K\). The number densities at all but the very earliest
times are too low to allow nuclei to be built up directly in many-body
collisions like \(2n+2p\to\mathrm{He}^4\). Complex nuclei must instead be
built up in sequences of two-body reactions, such as
\begin{equation}
  \begin{gathered}
    p+n\rightleftarrows d+\gamma,
    \\
    d+d\rightleftarrows\mathrm{He}^3+n
    \rightleftarrows\mathrm H^3+p,
    \\
    \mathrm H^3+d\rightleftarrows\mathrm{He}^4+n,
    \qquad\text{etc.}
  \end{gathered}
  \tag{15.7.32}\label{eq:15.7.32}
\end{equation}
There is no problem with the first step; the rate of deuterium production per
free neutron is
\begin{align}
  \lambda_d
  &=
  [4.55\times10^{-20}\,\mathrm{cm^3/sec}]\,n_p
  \notag\\
  &=
  27.4\,\mathrm{sec}^{-1}
  \left(
    \frac{a}{10^{-9}a_0}
  \right)^{-3}
  \left(
    \frac{\rho_{N0}}{10^{-30}\,\mathrm{g/cm^3}}
  \right)X_p,
  \tag{15.7.33}\label{eq:15.7.33}
\end{align}
and this is so much faster than the expansion rate \(1/t\) [see
Eq.~\eqref{eq:15.7.18}] that deuterons will appear with the equilibrium
abundance \eqref{eq:15.7.28}:
\begin{equation}
  X_d
  =
  \frac{3}{\sqrt2}X_pX_n\epsilon
  \exp\left(\frac{B_d}{kT}\right).
  \tag{15.7.34}\label{eq:15.7.34}
\end{equation}
However, no appreciable quantity of \(\mathrm H^3\), \(\mathrm{He}^3\),
\(\mathrm{He}^4\), or heavier nuclei can be formed until this equilibrium
deuterium abundance is high enough to allow \(d\)-\(d\), \(d\)-\(p\), or
\(d\)-\(n\) reactions to proceed at an adequate rate. According to
Table~\ref{tab:15.6}, the equilibrium deuterium abundance
\eqref{eq:15.7.34} is very small for \(T\) greater than about
\(0.8\times10^{9\,\circ}\mathrm K\). The low binding energy of deuterium
thus serves as a ``bottleneck,'' which delays the formation of complex nuclei
until \(T\) drops to near \(0.8\times10^{9\,\circ}\mathrm K\), or a little
earlier in models with a relatively high baryon number density.

Once nucleosynthesis begins, it proceeds very rapidly, because, according to
Table~\ref{tab:15.6}, any temperature less than
\(1.2\times10^{9\,\circ}\mathrm K\) is low enough to permit high equilibrium
concentrations of nuclei heavier than deuterons. However, it is not in fact
possible to produce appreciable quantities of elements heavier than helium
because, as mentioned above, the lack of stable nuclides with \(A=5\) or
\(A=8\) impedes nucleosynthesis via \(n\)-\(\alpha\),
\(p\)-\(\alpha\), or \(\alpha\)-\(\alpha\) collisions, whereas the Coulomb
barrier in the reactions
\(\mathrm{He}^4+\mathrm H^3\to\mathrm{Li}^7+\gamma\) and
\(\mathrm{He}^4+\mathrm{He}^3\to\mathrm{Be}^7+\gamma\) prevents them from
competing effectively with
\(p+\mathrm H^3\to\mathrm{He}^4+\gamma\) or
\(n+\mathrm{He}^3\to\mathrm{He}^4+\gamma\). Thus the effect of the nuclear
reactions \eqref{eq:15.7.32} is very rapidly to incorporate all available
neutrons into \(\mathrm{He}^4\) nuclei, which have by far the highest
binding energies of all nuclei with \(A<5\).

The nucleosynthesis process can only be followed in detail by a numerical
integration of a large number of rate equations. This has been done by
Peebles\hyperref[ch15-ref-93]{\textsuperscript{93}} for the reactions
\eqref{eq:15.7.32}, and by Wagoner, Fowler, and
Hoyle\hyperref[ch15-ref-110]{\textsuperscript{110}} for the reactions
\eqref{eq:15.7.32} plus the radiative processes
\begin{equation}
  \begin{gathered}
    p+d\rightleftarrows\mathrm{He}^3+\gamma,
    \qquad
    n+d\rightleftarrows\mathrm H^3+\gamma,
    \qquad
    p+\mathrm H^3\rightleftarrows\mathrm{He}^4+\gamma,
    \\
    n+\mathrm{He}^3\rightleftarrows\mathrm{He}^4+\gamma,
    \qquad
    d+d\rightleftarrows\mathrm{He}^4+\gamma,
  \end{gathered}
  \tag{15.7.35}\label{eq:15.7.35}
\end{equation}
plus a large number of other processes leading (weakly) up to nuclei as heavy
as \(\mathrm{Mg}^{24}\). Fortunately, none of these complications are
relevant to our basic problem, that of understanding the helium abundance.
All processes proceeding by strong and electromagnetic interactions, such as
the reactions \eqref{eq:15.7.32} and \eqref{eq:15.7.35}, will conserve the
total numbers of protons and neutrons. The only effect of nucleosynthesis on
the neutron--proton ratio is that, by ``turning off'' the decay of free
neutrons, it freezes this ratio at the value it had just before the onset of
nucleosynthesis. Before nucleosynthesis begins, the ratio of neutrons to all
nucleons is simply the quantity \(X_n\) given by Eq.~\eqref{eq:15.7.23}.
After nucleosynthesis is over, we have essentially nothing left but free
protons and \(\mathrm{He}^4\) nuclei, so the fraction of neutrons to all
nucleons is one-half the fraction of all nucleons that are bound in helium,
or one-half the abundance by weight of helium. Thus the abundance by weight of
cosmologically produced helium is simply given by
\begin{equation}
  Y
  =
  X_{\mathrm{He}^4}\;(\text{after nucleosynthesis})
  =
  2X_n\;(\text{just before nucleosynthesis}).
  \tag{15.7.36}\label{eq:15.7.36}
\end{equation}
According to the detailed calculations of Peebles, nucleosynthesis begins
abruptly at a temperature \(0.9\times10^{9\,\circ}\mathrm K\) for
\(\rho_{N0}=7\times10^{-31}\,\mathrm{g/cm^3}\), or at
\(1.1\times10^{9\,\circ}\mathrm K\) for
\(\rho_{N0}=1.8\times10^{-29}\,\mathrm{g/cm^3}\), just about as we should
expect from our qualitative considerations. Using Eq.~\eqref{eq:15.7.36}, we
can read off from Eq.~\eqref{eq:15.7.23} or Table~\ref{tab:15.5} that for
these two values of the present density, the helium abundance by weight
should be 26.2\% or 28.6\%, respectively. (Peebles%
\hyperref[ch15-ref-93]{\textsuperscript{93}} actually gives 25.8\% and
28.2\% in these two cases. This very slight discrepancy is simply due to the
small number of free neutrons that decay during the brief duration of
nucleosynthesis.) It is safe to say that in the class of cosmological models
considered here, a helium abundance by weight of about 27\% would be produced
cosmologically for any reasonable value of the present density. The reason
the helium abundance is so insensitive to the baryon number density is that
the neutron--proton ratio before nucleosynthesis is determined by the
interaction of the nucleons with the huge number of leptons, not with each
other, while the onset of nucleosynthesis is essentially determined by the
temperature, not the nucleon density.

Wagoner, Fowler, and Hoyle\hyperref[ch15-ref-110]{\textsuperscript{110}}
have calculated the cosmologically produced abundances, not only of the
isotopes of hydrogen and helium, but also of \(\mathrm{Li}^7\) and heavier
elements. Their results are given in Table~\ref{tab:15.7}. Note that the
abundances of all nuclides except \(\mathrm H^1\) and \(\mathrm{He}^4\) are
extremely small, so that production or destruction of these nuclides in stars
could have a serious effect on their observed ``cosmic'' abundances. For this
reason it is primarily the cosmic abundance of \(\mathrm{He}^4\) that serves
as a check on models of the early universe. (However, Geiss and Reeves%
\hyperref[ch15-ref-110a]{\textsuperscript{110a}} argue that the
\(\mathrm H^2\) and \(\mathrm{He}^3\) observed in the solar system do in fact
arise from the early universe. If this is correct, then the cosmic density
must be rather low, with a present value of order
\(3\times10^{-31}\,\mathrm{g/cm^3}\), to prevent the nuclear reactions,
which build up \(\mathrm H^2\) and \(\mathrm{He}^3\) into
\(\mathrm{He}^4\), from proceeding to completion.)

\begingroup
\renewcommand{\thetable}{15.7}
\begin{table}[htbp]
  \centering
  \caption{Cosmologically Produced Abundances (by Weight) of Various
  Nuclides, for Various Values of the Present Density \(\rho_{N0}\).}
  \label{tab:15.7}
  \scriptsize
  \setlength{\tabcolsep}{2.7pt}
  \resizebox{\linewidth}{!}{%
  \begin{tabular}{@{}lcccccccc@{}}
    \toprule
    &
    \multicolumn{8}{c}{\(\rho_{N0}\) (\(\mathrm{g/cm^3}\))} \\
    \cmidrule(lr){2-9}
    &
    \(10^{-31}\) &
    \(3.1\times10^{-31}\) &
    \(10^{-30}\) &
    \(3.1\times10^{-30}\) &
    \(10^{-29}\) &
    \(3.1\times10^{-29}\) &
    \(10^{-28}\) &
    \(3.1\times10^{-28}\) \\
    \midrule
    \(\mathrm H^1\) &
    0.763 & 0.748 & 0.737 & 0.728 & 0.719 & 0.709 & 0.701 & 0.691 \\
    \(\mathrm H^2\) &
    \(6.2\times10^{-4}\) &
    \(8.9\times10^{-5}\) &
    \(2.3\times10^{-5}\) &
    \(2.7\times10^{-7}\) &
    \(2.5\times10^{-12}\) &
    \(<10^{-12}\) &
    \(<10^{-12}\) &
    \(<10^{-12}\) \\
    \(\mathrm{He}^3\) &
    \(6.3\times10^{-5}\) &
    \(3.8\times10^{-5}\) &
    \(2.1\times10^{-5}\) &
    \(9.9\times10^{-6}\) &
    \(5.6\times10^{-6}\) &
    \(4.4\times10^{-6}\) &
    \(3.5\times10^{-6}\) &
    \(2.4\times10^{-6}\) \\
    \(\mathrm{He}^4\) &
    0.236 & 0.252 & 0.263 & 0.272 & 0.281 & 0.291 & 0.299 & 0.309 \\
    \(\mathrm{Li}^7\) &
    \(5.2\times10^{-10}\) &
    \(2.1\times10^{-10}\) &
    \(4.4\times10^{-9}\) &
    \(2.1\times10^{-8}\) &
    \(4.3\times10^{-8}\) &
    \(1.1\times10^{-7}\) &
    \(2.9\times10^{-7}\) &
    \(6.8\times10^{-7}\) \\
    Other &
    \(<10^{-12}\) &
    \(<10^{-12}\) &
    \(<10^{-12}\) &
    \(<10^{-12}\) &
    \(<10^{-12}\) &
    \(6\times10^{-12}\) &
    \(1.0\times10^{-10}\) &
    \(1.9\times10^{-9}\) \\
    \bottomrule
  \end{tabular}%
  }

  \vspace{0.5em}
  \begin{minipage}{0.96\linewidth}
    \footnotesize
    Results are taken from Tables 3A and 3B of R.~V. Wagoner, W.~A. Fowler,
    and F. Hoyle, \emph{Ap.\ J.} \textbf{148}, 3 (1967). (A value of
    \(3^\circ\mathrm K\) is assumed for the present black-body radiation
    temperature.)
  \end{minipage}
\end{table}
\endgroup

There are a number of different methods by which the helium abundance can be
measured in different parts of the universe.

\emph{(A) Stellar Masses and Luminosity.} The
theory\hyperref[ch15-ref-111]{\textsuperscript{111}} of stellar structure and
evolution allows us in principle (and even in practice) to calculate a star's
luminosity \(L\) as a function of time if we are given its mass \(M\) and
initial chemical composition. The chemical composition is usually specified
by three numbers, \(X,Y,Z\), defined as the fraction by weight respectively
of \(\mathrm H^1\), of \(\mathrm{He}^4\), and of everything else, with
\[
  X+Y+Z=1.
\]
(The heavy-element abundance \(Z\), though usually small, is an important
parameter for any star in radiative equilibrium, such as the sun, because it
determines the opacity of the star at a given density and temperature. The
helium abundance \(Y\) is important because it governs the mean molecular
weight appearing in the ideal-gas law.) If we can guess \(Z\) and the age of
a given star, then comparison of theory with measured values of \(M\) and
\(L\) allows us to compute \(Y\).

The best-studied star is, of course, the sun. Its mass and luminosity are
known quite accurately, and its age is believed to be close to the age of the
earth, or about \(4.5\times10^9\) years. From the absorption lines of hydrogen
and heavy elements it has been
estimated\hyperref[ch15-ref-112]{\textsuperscript{112}} that \(Z/X\) is
about 0.026 to 0.027 in the solar photosphere, though a more recent
study\hyperref[ch15-ref-113]{\textsuperscript{113}} gives
\(Z/X\simeq0.019\). (Unfortunately, helium lines are much too weak for
\(Y/X\) to be measured in the sun by this method.) Usually solar evolution
calculations are carried out for values of \(Z\) in the range 0.01 to 0.04.
At the time of the discovery of the cosmic microwave radiation the best solar
models\hyperref[ch15-ref-114]{\textsuperscript{114}} gave an initial helium
abundance \(Y=0.27\) for \(Z=0.02\) (or \(Y=0.32\) for \(Z=0.04\)), so it
was regarded as a great victory for the ``big bang'' cosmology that, with
\(T_{\gamma0}\simeq3^\circ\mathrm K\), it gives a primordial helium
abundance \(Y\simeq0.27\).

Unfortunately, this happy state lasted only until the advent of neutrino
astronomy. The same solar models that are used to calculate \(Y\) also can be
used to predict the flux of neutrinos from various nuclear reactions in the
sun. The sun derives its energy from the fusion of hydrogen into helium in a
proton--proton cycle, starting with the reactions
\[
  \begin{aligned}
    \mathrm H^1+\mathrm H^1
    &\to\mathrm H^2+e^++\nu
    &&(\overline E_\nu=0.263\,\mathrm{MeV}),
    \\
    \mathrm H^1+\mathrm H^1+e^-
    &\to\mathrm H^2+\nu
    &&(\overline E_\nu=1.4\,\mathrm{MeV}),
    \\
    \mathrm H^2+\mathrm H^1
    &\to\mathrm{He}^3+\gamma.
  \end{aligned}
\]
The cycle can then terminate with the ``PP I'' branch
\[
  \mathrm{He}^3+\mathrm{He}^3\to\mathrm{He}^4+2\mathrm H^1,
\]
or it can produce \(\mathrm{Be}^7\) by the reaction
\[
  \mathrm{He}^3+\mathrm{He}^4\to\mathrm{Be}^7+\gamma.
\]
In the latter case, one \(\mathrm{Be}^7\) nucleus and one proton are
converted into two \(\mathrm{He}^4\) nuclei by either the ``PP II'' branch
\[
  \begin{aligned}
    \mathrm{Be}^7+e^-
    &\to\mathrm{Li}^7+\nu
    &&(\overline E_\nu=0.80\,\mathrm{MeV}),
    \\
    \mathrm{Li}^7+\mathrm H^1
    &\to\mathrm{He}^4+\mathrm{He}^4,
  \end{aligned}
\]
or the ``PP III'' branch
\[
  \begin{aligned}
    \mathrm{Be}^7+\mathrm H^1
    &\to\mathrm B^8+\gamma,
    \\
    \mathrm B^8
    &\to\mathrm{Be}^8+e^++\nu
    &&(\overline E_\nu=7.2\,\mathrm{MeV}),
    \\
    \mathrm{Be}^8
    &\to2\mathrm{He}^4.
  \end{aligned}
\]
(Mean neutrino energies are given in parentheses.) Pontecorvo%
\hyperref[ch15-ref-115]{\textsuperscript{115}} and Alvarez%
\hyperref[ch15-ref-116]{\textsuperscript{116}} suggested that the neutrinos
could be detected in \(\mathrm{Cl}^{37}\) through the endothermic reaction
\begin{equation}
  \nu+\mathrm{Cl}^{37}\to e^-+\mathrm{Ar}^{37}.
  \tag{15.7.37}\label{eq:15.7.37}
\end{equation}
The \(\mathrm{Ar}^{37}\) decays by electron capture with a convenient
half-life of 35 days, so it can be detected by its radioactivity after
chemical separation. As pointed out by Bahcall,%
\hyperref[ch15-ref-117]{\textsuperscript{117}} the energetic neutrinos from
\(\mathrm B^8\) beta decay are particularly effective in the reaction
\eqref{eq:15.7.37}, because they can induce superallowed transitions to an
excited state of \(\mathrm{Ar}^{37}\). Hence, even though the PP III branch
is much less important than the PP II branch, about 90\% of the neutrino
absorption events in \(\mathrm{Cl}^{37}\) would be expected to arise from
\(\mathrm B^8\) neutrinos, and about 10\% from \(\mathrm{Be}^7\) neutrinos.
Using the extant solar models\hyperref[ch15-ref-114]{\textsuperscript{114}}
with \(Y=0.27\), Bahcall\hyperref[ch15-ref-117]{\textsuperscript{117}}
calculated a neutrino capture rate on the earth of
\((4\pm2)\times10^{-35}\,\mathrm{sec}^{-1}\) per \(\mathrm{Cl}^{37}\) atom,
and Davis\hyperref[ch15-ref-118]{\textsuperscript{118}} set out to measure
this rate, using 100,000 gallons of perchlorethylene
\((\mathrm C_2\mathrm{Cl}_4)\), a common cleaning fluid, in the Homestake
gold mine at Lead, South Dakota. In 1968 Davis et al.%
\hyperref[ch15-ref-119]{\textsuperscript{119}} announced that they had
failed to detect any solar neutrinos and could set an upper limit of
\(0.3\times10^{-35}\,\mathrm{sec}^{-1}\) on the absorption rate per
\(\mathrm{Cl}^{37}\) atom, about an order of magnitude less than what had
originally been expected! This discrepancy between theory and observation,
in the first experiment that ever looked directly into the solar interior,
has shaken the general faith in accepted solar models, and in the values they
yield for the initial helium abundance of the sun. Needless to say, a great
deal of work has been put into recalculation of the expected neutrino fluxes,
using improved values for the opacity and for various nuclear reaction rates.
In a companion paper to the letter of Davis et al.,%
\hyperref[ch15-ref-119]{\textsuperscript{119}} Bahcall et al.%
\hyperref[ch15-ref-120]{\textsuperscript{120}} estimated an absorption rate
for \(Z=0.015\) of
\((0.75\pm0.3)\times10^{-35}\,\mathrm{sec}^{-1}\) per
\(\mathrm{Cl}^{37}\) atom, still too large by a factor of two, and
calculations using the Berkeley stellar structure code gave slightly larger
rates.\hyperref[ch15-ref-121]{\textsuperscript{121}} Iben%
\hyperref[ch15-ref-122]{\textsuperscript{122}} has noted that both \(Y\) and
the neutrino absorption rates are increasing functions of \(Z\), and that the
minimum possible absorption rate, attained by \(Z=0\) and \(Y\simeq0.17\),
is just about the same as Davis et al.'s upper limit. The latest calculation
of Bahcall and Ulrich\hyperref[ch15-ref-122a]{\textsuperscript{122a}} gives
a counting rate of
\((0.9\pm0.5)\times10^{-35}\,\mathrm{sec}^{-1}\) per
\(\mathrm{Cl}^{37}\) atom.

Meanwhile, Davis' group continued their observations, and have recently
announced a counting rate of
\((0.15\pm0.1)\times10^{-35}\,\mathrm{sec}^{-1}\) per
\(\mathrm{Cl}^{37}\) atom,\hyperref[ch15-ref-122b]{\textsuperscript{122b}}
about six times less than expected. In view of this discrepancy, the question
of the initial solar helium abundance must for the present be regarded as
unsettled.

The masses and luminosities are also known for a number of nearby Population
I stars that happen to belong to binary systems. Comparison of these \(M\)
and \(L\) values with the theoretical \(Y\)-dependent \(M\)--\(L\) relation
yields \(Y\) values\hyperref[ch15-ref-123]{\textsuperscript{123}} for these
stars lying generally in the range 0.25 to 0.35. It would be very interesting
to carry out this analysis for stars of Population II, because they represent
an earlier stellar generation, and also because the Davis neutrino experiment
has shaken our faith in the theory of stars of Population I. Unfortunately,
there are very few stars of Population II near the sun. One of them,
\(\mu\) Cassiopeiae A, belongs to a binary system whose separation has
recently been measured by a very ingenious
method.\hyperref[ch15-ref-124]{\textsuperscript{124}} The resulting mass
value, together with the observed values for \(L\) and \(Z/X\), does not
agree with theory\hyperref[ch15-ref-125]{\textsuperscript{125}} for any
value of \(Y\), but fits best a low helium abundance with \(Y\lesssim0.05\).
However, the validity of this mass determination has since been put into
doubt.\hyperref[ch15-ref-125a]{\textsuperscript{125a}}

\emph{(B) Direct Solar Measurements.} There are a number of different
methods for estimating the present helium abundance of the sun, which do not
rest on any detailed theory of solar structure and evolution. Measurements
of the ratio \(Y/Z\) in solar cosmic
rays,\hyperref[ch15-ref-126]{\textsuperscript{126}} together with the
spectroscopic determination of \(Z/X\) in the solar photosphere mentioned
above, suggest a helium
abundance\hyperref[ch15-ref-127]{\textsuperscript{127}}
\(Y\simeq0.20\) to 0.26. During periods of quiet sun, the
\(\mathrm{He}^4/\mathrm H\) ratio in the solar wind suggests a
value\hyperref[ch15-ref-128]{\textsuperscript{128}} \(Y\) about equal to
0.15, but the helium content of the solar wind roughly doubles during
magnetic storms.\hyperref[ch15-ref-129]{\textsuperscript{129}}
Unfortunately, the solar surface is too cool to allow a spectroscopic
determination of \(Y\), but a value of \(Y\simeq0.38\) is suggested by
observation of solar prominences.\hyperref[ch15-ref-130]{\textsuperscript{130}}

\emph{(C) Globular Clusters: Theory.} As already mentioned in
Section~\ref{sec:15.3}, the comparison of the numbers of stars in different
regions of the Hertzsprung--Russell diagrams of globular clusters with theory
yields results for both the age and the initial helium abundance of these
clusters. Iben\hyperref[ch15-ref-131]{\textsuperscript{131}} derives values
of \(Y\) in the range of 0.24 to 0.33, corresponding to ages
\(18\times10^9\) years to \(9\times10^9\) years. Comparison of the stellar
pulsation theory of Christy\hyperref[ch15-ref-132]{\textsuperscript{132}}
with the location of variable stars in the Hertzsprung--Russell diagrams of
the globular clusters M3, M15, and M92 yields \(Y\simeq0.26\) to 0.32 for
these clusters.\hyperref[ch15-ref-133]{\textsuperscript{133}} These studies
should be given particular weight, because the globular clusters are believed
to be among the first objects to condense out of the primordial gas of
hydrogen and helium.

\emph{(D) Stellar Spectra.} Helium lines are visible in the photospheres of a
large number of hot stars of both populations. In general, helium abundances
appear to be high,\hyperref[ch15-ref-123]{\textsuperscript{123}} with \(Y/X\)
of order 0.4, and some stars are apparently superabundant in helium. There are
several classes of old stars that show anomalously weak helium lines, such as
the blue Population II stars on the horizontal branch of the
Hertzsprung--Russell diagram of globular
clusters.\hyperref[ch15-ref-134]{\textsuperscript{134}} One peculiar star,
3 Centauri A, has a low abundance of helium, most of which is in the form of
the isotope \(\mathrm{He}^3\)! Planetary
nebulae\hyperref[ch15-ref-135]{\textsuperscript{135}} and
novae\hyperref[ch15-ref-123]{\textsuperscript{123}} generally appear to be
superabundant in helium.

\emph{(E) Spectroscopy of Interstellar Matter.} Optical-frequency emission
lines from H II regions in our galaxy yield helium--hydrogen number
ratios\hyperref[ch15-ref-123]{\textsuperscript{123}} that are consistently
in the range 0.10--0.14, corresponding to a helium abundance by weight
\(Y\simeq0.27\)--0.36. It is also possible to observe the recombination of
ionized helium at radio frequencies;%
\hyperref[ch15-ref-136]{\textsuperscript{136}} the radiation emitted in a
transition \(n+1\to n\) has a wavelength proportional to \(n^3\) for
\(n\gg1\), so that transitions with \(n\simeq100\) have wavelengths of the
order of centimeters. The helium--hydrogen number
ratios\hyperref[ch15-ref-123]{\textsuperscript{123}} deduced from radio
observations of interstellar matter range from 0.06 to 0.16, corresponding to
\(Y\simeq0.14\) to \(Y\simeq0.40\).

\emph{(F) Extragalactic Measurements.} The emission lines of helium
observed\hyperref[ch15-ref-123]{\textsuperscript{123}} in H II regions in
galaxies within and without our local group indicate a helium abundance
similar to that of the H II regions in our own galaxy. On the other hand,
quasi-stellar sources show remarkably weak helium
lines.\hyperref[ch15-ref-123]{\textsuperscript{123}}

There is clearly a good deal of evidence for a universal helium abundance by
weight not too different from the predicted value of 27\%. Unfortunately,
there is also a large body of evidence for a much smaller helium abundance.
The clarification of this problem would be of the highest importance for
cosmology, because the cosmologically produced helium, together with the
\(2.7^\circ\mathrm K\) radiation background, may be the only relics of the
primordial fireball that can serve as clues to the early history of the
universe.

In order to keep an open mind about the synthesis of the elements in the
early universe, it is useful to consider the possible modifications in
physical or astrophysical theory that could affect the production of helium:

\emph{(A) Cool Models.} If the observed microwave background proves not to
be black-body radiation left over from the early universe, then we would have
to face the possibility that the true present black-body temperature
\(T_{\gamma0}\) is very much less than \(2.7^\circ\mathrm K\). In this case,
the baryon number density at any given past temperature could be much greater
than assumed above, with a consequent increase in the rate of nuclear
reactions and in the abundance of complex nuclei produced in the early
universe. Indeed, it was the high helium abundance produced in such cool
models that led Gamow and his
collaborators\hyperref[ch15-ref-51]{\textsuperscript{51}} to suggest the
presence of a hot radiation background.

\emph{(B) Fast or Slow Models.} Various mechanisms might increase or
decrease the expansion rate. In particular, if the universe contained a
thermal distribution of additional massless quanta, such as gravitons,
Brans--Dicke scalar particles, or new kinds of neutrinos, then the energy
density at a given temperature would be greater, and so, according to
Eq.~\eqref{eq:15.6.44}, the time required to reach that temperature would be
shorter. The rate \eqref{eq:15.7.33} of deuteron production per free neutron
is normally larger than the expansion rate at
\(T=10^{9\,\circ}\mathrm K\) by a factor 10 to \(10^3\) (for a present
density of \(10^{-31}\,\mathrm{g/cm^3}\) to
\(10^{-29}\,\mathrm{g/cm^3}\)), so for a moderate shortening of the time
scale there would still be plenty of time for nucleosynthesis to occur at
\(T\simeq10^{9\,\circ}\mathrm K\). In this case, the only effect of the
faster expansion would be to cut down the time available for the conversion
of neutrons into protons, so that the neutron fraction at
\(10^{9\,\circ}\mathrm K\) would be closer to its initial value of
\(\tfrac12\), and more helium would be produced. However, if the time scale
were extremely short, there would not be time for complex nuclei to be formed
before the density (and, for \(\mathrm{He}^3\) and \(\mathrm{He}^4\)
formation, the temperature) falls too low. The detailed calculations of
Peebles\hyperref[ch15-ref-137]{\textsuperscript{137}} show that for
\(T_{\gamma0}=3^\circ\mathrm K\) and a present density of
\(7\times10^{-31}\,\mathrm{g/cm^3}\) to
\(1.8\times10^{-29}\,\mathrm{g/cm^3}\), the \(\mathrm{He}^4\) abundance
rises as the time scale is shortened until it reaches a maximum of 60\% to
80\% (by weight) for a time scale shortened by a factor \(10^{-1}\) to
\(10^{-2}\), and then falls off again. The deuteron abundance continues to
rise with shortening time scale, reaching a maximum of about 9\% (by weight)
when the time scale is shortened by a factor \(3\times10^{-3}\) to
\(3\times10^{-4}\), and then falls off again. On the other hand, if the
expansion time scale were somehow lengthened, the only effect would be that
more neutrons would decay into protons before nucleosynthesis occurs, so that
less helium would be produced.

\emph{(C) Neutrino--Electron Interactions.} The thermal history of the early
universe was worked out in the last section under the assumption that both
electron- and muon-type neutrinos lose thermal contact with the
\(e^+-e^--\gamma\) plasma before the onset of \(e^+-e^-\) annihilation.
This assumption is probably valid if the neutrino--electron scattering is
produced by the usual Fermi weak interaction with the same strength as in
nuclear beta decay or muon decay. However, the neutrino--electron interaction
has not yet been measured experimentally, and it could be somewhat stronger
than expected.\hyperref[ch15-ref-138]{\textsuperscript{138}} In this case,
the \(\nu_e\) and \(\bar\nu_e\) (and possibly also the \(\nu_\mu\) and
\(\bar\nu_\mu\)) might remain in thermal equilibrium with the
\(e^+-e^--\gamma\) plasma until nearly all \(e^+-e^-\) pairs have
annihilated. The effect would be to increase the energy density at any given
temperature, and also to eliminate the distinction between \(T_\nu\) and
\(T\) in the rates \(\lambda(n\to p)\) and \(\lambda(p\to n)\). Detailed
calculations\hyperref[ch15-ref-139]{\textsuperscript{139}} show that if
electron-type neutrinos remain in thermal equilibrium until helium synthesis,
then the abundance of cosmologically produced helium would be about 29\%
instead of 27\%.

\emph{(D) Neutrino or Antineutrino Degeneracy.} It is also interesting to
consider the effect of a \(\nu_e\) or \(\bar\nu_e\) degeneracy on the helium
abundance. One effect is that the increased density would speed up the
expansion. In addition, the imbalance between neutrinos and antineutrinos
would affect the relative abundance of protons and neutrons. The difference
between the neutron and proton chemical potentials in thermal equilibrium is
given by Eq.~\eqref{eq:15.6.4}:
\[
  \mu_n-\mu_p=\mu_{e^-}-\mu_{\nu_e}.
\]
We saw in the last section that, during the period of interest,
\(\mu_{e^-}\) is required to be negligible to maintain charge neutrality,
whereas \(\mu_{\nu_e}/kT\) is a constant \(\nu\) (with
\(\lvert\nu\rvert\lesssim45\)):
\[
  \mu_{e^-}\simeq0,
  \qquad
  \mu_{\nu_e}\simeq\nu kT.
\]
The equilibrium neutron fraction is then given by
Eq.~\eqref{eq:15.7.19} as
\[
  X_n
  \equiv
  \frac{n_n}{n_p+n_n}
  =
  \left[
    1+\exp\left(\nu+\frac{Q}{kT}\right)
  \right]^{-1},
\]
where \(Q\equiv m_n-m_p\). Thus, if \(\nu\) were large and positive, the
neutron fraction would start small and remain small, so that very little
nucleosynthesis would occur. If \(\nu\) were moderately negative, say
\(\nu\simeq-1\), then the initial neutron fraction would be high, so that
after some neutrons were converted to protons, the neutron fraction at the
onset of nucleosynthesis could be close to the optimum value of 50\%, and
essentially all the matter of the universe could be converted to helium. If
\(\nu\) were large and negative, then the initial neutron abundance would be
extremely high, and no nucleosynthesis could occur until some neutrons could
decay into protons, at which time the nucleon density would have been too low
to allow much synthesis of complex nuclei. Detailed calculations of the
abundance of \(\mathrm H^2,\mathrm{He}^3,\mathrm{He}^4\), and
\(\mathrm{Li}^7\) as functions of \(\nu\) have been carried out by Wagoner,
Fowler, and Hoyle,\hyperref[ch15-ref-140]{\textsuperscript{140}} taking into
account the effects of the neutrino or antineutrino degeneracy on the rates
\eqref{eq:15.7.2}--\eqref{eq:15.7.7}. These calculations show that if the
``missing mass'' consists of degenerate neutrinos or antineutrinos with
\(\lvert\nu\rvert\simeq30\), then the cosmologically produced abundance
(by weight) of helium would be considerably less than 1\%. On the other hand,
if the lepton number density \(N_E\) of the universe is of the same order as
the baryon number density \(N_B\), then \eqref{eq:15.6.52} shows that
\(\lvert\nu\rvert\) is of order \(1/\sigma\), or about \(10^{-9}\). [See
Eq.~\eqref{eq:15.5.15}.] In this case the slight excess of neutrinos or
antineutrinos has no appreciable effect on the synthesis of helium.

One final warning: Even if a large cosmic abundance of helium is definitely
established, it would not necessarily follow that this helium was formed in
the early universe. Geoffrey
Burbidge\hyperref[ch15-ref-141]{\textsuperscript{141}} has particularly
emphasized the possibility that helium could have been synthesized in an
earlier, more luminous phase of our galaxy's history, perhaps in massive
galactic objects. A good part of the calculations discussed in this chapter
would also apply to nucleosynthesis in the collapse of massive
stars.\hyperref[ch15-ref-142]{\textsuperscript{142}}
